\subsection{Cubesets: the Basics}\label{ssec:cubesets-basics}

Now that we have established what the basic building blocks of our geometric objects are,
namely the $n$-dimensional cubes ${\bbox}^n$, we can start creating geometric spaces out of them.
As usual, one wants to glue multiple cubes together to obtain \emph{cubesets}.
This is achieved by taking the free cocompletion of the category $\bbox$.

\subsubsection{Cubesets are Presheaves}

\begin{definition}\label{def:cubeset}
The category of cubesets $\cSet$ is the presheaf category on $\bbox$,
\[
    \cSet = \PSh(\bbox) = \Fun({\bbox}^{\op},\Set).
\]
\end{definition}

As a presheaf category, the category of cubesets $\cSet$ enjoys many pleasant properties.
One of the most important certainly is the fact that it is a Grothendieck topos
where all limits and colimits can be computed pointwise.

Recall that the Yoneda embedding $\bbox\to\PSh(\bbox)$ exhibits $\cSet$ as the free cocompletion
of the cube category $\bbox$.
This is described by the following universal property of the category $\cSet$ (by which it is uniquely determined):
every functor $\bbox\to\mcD$ where $\mcD$ has all small colimits
has a unique extension to a colimit preserving functor $\cSet\to\mcD$.
Concretely, this means that every presheaf $F$ can be represented as the canonical colimit over its category of elements,
or equivalently over the slice category $\bbox_{/F}$, that is over all cubes ${\bbox}^n\to F$ in $F$:
\[
    F = \varinjlim_{{\bbox}^n\to F}{\bbox}^n
\]
Moreover, this colimit is \emph{free} in the sense that for
$F = \varinjlim_{{\bbox}^n\to F} {\bbox}^n$ and $G = \varinjlim_{{\bbox}^m\to G} {\bbox}^m$,
the double limit formula
\[
    \hom(F,G) = \varprojlim_{{\bbox}^n\to F}\varinjlim_{{\bbox}^m\to G}\hom_{\bbox}({\bbox}^n,{\bbox}^m)
\]
holds.
This means that every cube ${\bbox}^n$ in $\cSet$ is a tiny object and in particular compact projective.
Thus, the category $\PSh(\bbox)$ is $\aleph_0$-presentable, so every colimit preserving functor
$G\colon\PSh(\bbox)\to\mcD$ admits a right adjoint, given on an object $D\in\mcD$ by the presheaf
\[
    \langle n\rangle\mapsto\hom(G({\bbox}^n),D).
\]
The necessity of defining the right adjoint as it is given follows from the adjunction and the Yoneda lemma.
Moreover, if $G$ preserves limits and is accessible
(i.e. preserves $\lambda$-filtered colimits for $\lambda$ big enough), and $\mcD$ is presentable,
then it admits a left adjoint (and conversely).


Another important construction is the lift of the symmetric monoidal structure
${\bbox}^n\otimes{\bbox}^m={\bbox}^{n+m}$ in $\bbox$ onto the full $\cSet$.
This is achieved by \emph{Day convolution}.

\begin{remark}[Day Convolution]
If $\mcC$ is a monoidal category, then one can define a tensor product on its presheaf category $\PSh(\mcC)$.
Given $A,B\in\PSh(\mcC)$, their tensor product $A\otimes B$ is given by left Kan extension
of the external product $A\overline{\times}B\colon\mcC^{\op}\times\mcC^{\op}\to\set$ along the tensor product
$\otimes^\op\colon\mcC^\op\times\mcC^\op\to\mcC$,
\begin{center}
\begin{tikzcd}
	{\mcC^\op\times\mcC^\op} && \mcC^\op \\
	{\set\times\set} \\
	\set
	\arrow["\otimes^\op", from=1-1, to=1-3]
	\arrow["{A\times B}"', from=1-1, to=2-1]
	\arrow[dashed, from=1-3, to=3-1]
	\arrow["\times"', from=2-1, to=3-1]
\end{tikzcd}
\end{center}
which is given by the coend
\[
    (A\otimes B)(e) = \int^{(c,d)\in\mcC\times\mcC}\hom(e,c\otimes_{\mcC}d)\times A(c)\times B(d)
\]
or equivalently, the colimit
\[
    (A\otimes B)(e) = \varinjlim_{e\to c\otimes d} (A(c)\times B(d)).
\]
This defines a monoidal structure on $\PSh(\mcC)$.
If the tensor product on $\mcC$ is symmetric, then so is the Day convolution product.
\end{remark}

Specialized to cubesets it is surprising to see that Day convolution agrees with the cartesian monoidal structure.

\begin{proposition}\label{prop:Day-conv-cubeset}
Day convolution on $\cSet = \PSh(\bbox)$ agrees with the cartesian monoidal structure.
\end{proposition}

\begin{proof}
Recall that on $\bbox$, the symmetric monoidal structure defined by ${\bbox}^n\otimes{\bbox}^m = {\bbox}^{n+m}$ is the cartesian one, see Proposition~\ref{prop:bbox-symm-monoid}.
Then the general case follows from the Kan formula and the fact that colimits are stable under base change in $\PSh(\bbox)$.
\end{proof}

\begin{remark}
The same line of reasoning also works when replacing $\bbox$ by $\bbGamma$:
the proof of the preceding Proposition~\ref{prop:Day-conv-cubeset} relies only on the fact
that on the site $\mcC$, the monoidal structure is the cartesian one.
But this is also the case for $\bbGamma = \bblox$ by Remark~\ref{rem:bblox-symm-monoid}.
\end{remark}

Next, let us discuss the internal Hom in $\cSet$.
That Day convolution agrees with the cartesian product also implies that the tensor product is closed.
Recall that in any category of presheaves, the internal Hom of two cubesets $X$ and $Y$ is given on objects by
\[
    \ihom(X,Y)(n) = \hom(X\times{\bbox}^n,Y).
\]
This description makes $\cSet$ into a closed symmetric monoidal category where the tensor product
is the cartesian product (or equivalently, Day convolution).

\begin{example}
The \emph{path space} $\uP X$ of a cubeset $X$ is given by $\uP X = \ihom([1],X)$,
which has simply $\uP X(n) = X(n+1)$, i.e., the cubeset of edges in $X$.

\end{example}

\begin{example}\label{ex:host_kra_cubegroups}
An important class of examples of cubesets is given by \emph{Host-Kra cubegroups}.
They are defined as follows:
For a filtered group $G_\bullet=(G_0=G_1\ge G_2\ge \cdots)$ with $[G_i,G_j]\subset G_{i+j}$
define $\HK(G_\bullet)$ to be the sub-cubeset of $i_*(G_0)$
whose $n$-cubes are the subgroup of $G^{2^n}$ spanned by configurations
\[
    g^F(\omega)=\begin{cases}
        g,\quad & \omega\in F\\
        1,\quad & \omega \notin F,
    \end{cases},\qquad \omega\in \{0,1\}^n
\]
for a face $F$ of $\{0,1\}^n$ and $g\in G_{\mathrm{codim} F}$.

In this paper, the most relevant case of these cubegroups is given by an abelian group $A$
with filtration $A=G_0=\dots=G_k\ge 0=G_{k+1}$.
The corresponding cubegroup is called the \emph{$k$-th Eilenberg-Mac Lane cubeset $\cD_k(A)$}.
For more on cubegroups we refer to \cite{Bihlmaier2026b}.
\end{example}

\begin{remark}\label{rem:EM_preserves_products}
Note that the Eilenberg-Mac Lane cubeset induces a functor $\cD_k\colon \Ab\to \cSet$
for every $k\geq 1$ that factors through \emph{abelian} cubegroups.
This functor preserves products.
\end{remark}


\subsubsection{Cube- and \texorpdfstring{$\bbGamma$-}{Gamma}sets}

Since every morphism of sites (the indexing categories) induces a triple of adjunctions
between the corresponding presheaf categories, and $\bbGamma$ is a retract of the unstraightening $\bbox$,
we obtain several comparison functors between cubesets and $\bbGamma$-sets.

Let us first consider the triple of adjunctions induced by the unstraightening $\mathrm{Un}\colon \bbox\to \bbGamma$.

\begin{remark}
The cocartesian fibration/Grothendieck opfibration $\mathrm{Un}\colon\bbox\to\bbGamma$ induces a triple of adjunctions
\begin{center}
\begin{tikzcd}
	{\PSh(\bbox)} && {\PSh(\bbGamma)}
	\arrow[""{name=0, anchor=center, inner sep=0}, "{{\mathrm{Un}_\ast}}"', curve={height=24pt}, from=1-1, to=1-3]
	\arrow[""{name=0p, anchor=center, inner sep=0}, phantom, from=1-1, to=1-3, start anchor=center, end anchor=center, curve={height=24pt}]
	\arrow[""{name=1, anchor=center, inner sep=0}, "{{\mathrm{Un}_!}}", curve={height=-24pt}, from=1-1, to=1-3]
	\arrow[""{name=1p, anchor=center, inner sep=0}, phantom, from=1-1, to=1-3, start anchor=center, end anchor=center, curve={height=-24pt}]
	\arrow[""{name=2, anchor=center, inner sep=0}, "{{\mathrm{Un}^\ast}}"', from=1-3, to=1-1]
	\arrow[""{name=2p, anchor=center, inner sep=0}, phantom, from=1-3, to=1-1, start anchor=center, end anchor=center]
	\arrow[""{name=2p, anchor=center, inner sep=0}, phantom, from=1-3, to=1-1, start anchor=center, end anchor=center]
	\arrow["\dashv"{anchor=center, rotate=-81}, draw=none, from=1p, to=2p]
	\arrow["\dashv"{anchor=center, rotate=-99}, draw=none, from=2p, to=0p]
\end{tikzcd}
\end{center}
where $\mathrm{Un}_!$ and $\mathrm{Un}_\ast$ are the pointwise left, resp. right, Kan extension
of $\mathrm{Un}$ along the Yoneda embedding $\bbox\to\PSh(\bbox)$.
The right adjoint $\mathrm{Un}_\ast$ is given by
\[
    \mathrm{Un}_\ast X(n) = \hom_{\PSh(\bbGamma)}(\langle n\rangle,\mathrm{Un}_\ast X) = \hom_{\cSet}(\mathrm{Un}^\ast\langle n\rangle,X)
\]
where $\mathrm{Un}^\ast\langle m\rangle = y(\langle m\rangle)\circ\mathrm{Un}$.
The left adjoint can be computed via the formula
\begin{align*}
    \mathrm{Un}_!X(m) &= \varinjlim_{\bbox^n\to\langle m\rangle}X(n),
\end{align*}
where the index category runs over all maps ${\bbox}^n\to\langle m\rangle$ for fixed $m$.
It equivalently can be described as the unique cocontinuous extension of the horizontal functor
\begin{center}
\begin{tikzcd}
	\bbox & \bbGamma & {\PSh(\bbGamma)} \\
	\\
	\cSet
	\arrow["{\mathrm{Un}}", from=1-1, to=1-2]
	\arrow["y"', from=1-1, to=3-1]
	\arrow["y", from=1-2, to=1-3]
	\arrow["{\mathrm{Un}_!}", dashed, from=1-3, to=3-1]
\end{tikzcd}
\end{center}
making the diagram commutative, as in the universal property of the free cocompletion.
Indeed, for any $X\in\PSh(\bbGamma)$, we compute that
\begin{align*}
    \hom_{\PSh(\bbGamma)}(\mathrm{Un}_! y_{\cSet}(n),X) 
    &= \hom_{\cSet}(y_{\cSet}(n),\mathrm{Un}^\ast X) \\
    &= \mathrm{Un}^\ast X(n) \\
    &= X(\mathrm{Un}(n)) \\
    &= X(n)  \\
    &= \hom_{\PSh(\bbGamma)}(\langle n\rangle,X)
\end{align*}
and so the claim follows from the Yoneda lemma.

Such a triple of adjoint functors always induces a geometric morphism $\mathrm{Un}\colon\cSet\to\PSh(\bbGamma)$
whose right adjoint part is $\mathrm{Un}_\ast$ and the left exact left adjoint is $\mathrm{Un}^\ast$.
\end{remark}

Another important example is that of the wide subcategory inclusion $j\colon\bbGamma=\bblox\hookrightarrow\bbox$, which as before induces a triple of adjoint functors.

\begin{remark}\label{rem:adjoint-triple-j}
Consider the triple of adjunctions
\begin{center}
\begin{tikzcd}
	\cSet && {\PSh(\bbGamma)}
	\arrow[""{name=0, anchor=center, inner sep=0}, "{{j^\ast}}", from=1-1, to=1-3]
	\arrow[draw=none, from=1-1, to=1-3]
	\arrow[""{name=1, anchor=center, inner sep=0}, "{{j_\ast}}", curve={height=-24pt}, from=1-3, to=1-1]
	\arrow[""{name=2, anchor=center, inner sep=0}, "{{j_!}}"', curve={height=24pt}, from=1-3, to=1-1]
	\arrow["\dashv"{anchor=center, rotate=-65}, draw=none, from=0, to=1]
	\arrow["\dashv"{anchor=center, rotate=-115}, draw=none, from=2, to=0]
\end{tikzcd}
\end{center}
Here, $j^*$ is the pullback of $j$, simply forgetting the affine structure of a cubeset and only remembering the linear maps.
The left adjoint $j_!$ is the unique cocontinuous extension of the functor $\bbGamma\to\cSet$ assigning $\langle n\rangle$ to ${\bbox}^n$.
We have that $j_!(y_{\bbGamma}(n)) = y_{\bbox}(n)$.
Indeed, for any $X\in\cSet$, we compute that
\[
    \hom_{\cSet}(j_! y_{\bbGamma}(n),X) = \hom_{\PSh(\bbGamma)}(y_{\bbGamma}(n),j^\ast X) = j^\ast X(n) = X(j(n)) = X(n) = \hom_{\cSet}({\bbox}^n,X)
\]
and so the claim follows from the Yoneda lemma.
\end{remark}

More interestingly, in this case not only is $j_\ast\colon\PSh(\bbGamma)\to\cSet$
(the right adjoint part of) a geometric morphism but also $j^\ast\colon\cSet\to\PSh(\bbGamma)$.

\begin{lemma}\label{lem:j-geometric-morphism}
The functor $j_!$ preserves finite limits and finite unions of subobjects.
\end{lemma}

\begin{proof}
For $X$ a presheaf on $\bbGamma$, by the formula for the pointwise Kan extension we have that
\[
    j_! X({\bbox}^n) = \varinjlim_{{\bbox}^n\to j(\langle m\rangle)}X(\langle m\rangle)
\]
where we denote by ${\bbox}^n$ an object of $\bbox$ and by $\langle m\rangle$ an object in $\bbGamma$,
and the index diagram is the projection
\[
    Q\colon ({\bbox}^n\downarrow j)\to\bbGamma
\]
where $({\bbox}^n\downarrow j)$ is the category of objects in $\bbGamma$ $j$-over ${\bbox}^n$.
By Lemma~\ref{lem:fin-star-co-comp}, $\bbGamma$ has all finite limits
and by Lemma~\ref{lem:bblox-bbox-pres-limits}, $j$ preserves them.
Now it is a general fact that if $\bbGamma$ is finitely complete and $j$ preserves finite limits,
then also $({\bbox}^n\downarrow j)$ has finite limits and the projection $Q$ creates them.
This means that the index category in the colimit formula for $j_! X$ is cofiltered,
which in turn means that, since finite limits commute with filtered colimits in $\set$,
that $j_!$ preserves finite limits.

Being a left adjoint, $j_!$ preserves colimits.

The union of two subobjects is given by the pushout of their pullback,
and $j_!$ preserves both those operations, so it preserves binary unions of subobjects.
By induction, it preserves all finite unions of subobjects.
\end{proof}

\begin{remark}
For an example of a $\bbGamma$-set which is not a cubeset (doesn't arise from $j^\ast X$ for some $X\in\cSet$)
consider the \emph{upper Host-Kra cubegroups}.
For a filtered group $G_\bullet$ the upper Host-Kra $\bbGamma$-group $\HK_*(G_\bullet)$ can be constructed
by defining the value on $\langle n\rangle $ to be the subgroup of $G^{\{0,1\}^n}$ generated by all elements of the form
\[
    g^F\colon\{0,1\}^n\to G,\quad \omega\mapsto \begin{cases}
        g,\quad & \text{if}\;\omega\notin F, \\
        e,\quad & \text{else},
    \end{cases}
\]
for $F\subsetneqq \{0,1\}^n$ being an inert face of codimension $i>0$ and $g\in G_i$.
Compare also with the full Host-Kra cubegroup~\ref{ex:host_kra_cubegroups},
see also \cite{Bihlmaier2026b}.
\end{remark}

A further nice application of adjoint functors induced by morphisms of sites is the comparison to plain sets, and its adjoints.

\begin{lemma}\label{lem:point-cset}
Consider the embedding of the terminal object $i\colon\ast\to\bbox$.
This induces a triple of adjoints
\begin{center}
\begin{tikzcd}
	\cSet && {\Set.}
	\arrow[""{name=0, anchor=center, inner sep=0}, "{{i^\ast}}", from=1-1, to=1-3]
	\arrow[draw=none, from=1-1, to=1-3]
	\arrow[""{name=1, anchor=center, inner sep=0}, "{{i_\ast}}", curve={height=-24pt}, from=1-3, to=1-1]
	\arrow[""{name=2, anchor=center, inner sep=0}, "{{i_!}}"', curve={height=24pt}, from=1-3, to=1-1]
	\arrow["\dashv"{anchor=center, rotate=-65}, draw=none, from=0, to=1]
	\arrow["\dashv"{anchor=center, rotate=-115}, draw=none, from=2, to=0]
\end{tikzcd}
\end{center}
Here, the pullback $i^\ast\colon\PSh(\bbox)\to \Set=\PSh(\ast)$ of $i\colon\ast\to\bbox$ is the evaluation at $0$.
Its left adjoint $i_!$ sends a set $S$ to the constant presheaf ${\bbox}^n\mapsto S$.
The right adjoint $i_\ast$ assigns to a set $S$ the cubeset given by
\[
    {\bbox}^n\mapsto\hom_{\set}(\{0,1\}^n,S) =  S^{2^n}.
\]
Both adjoints $i_!$ and $i_\ast$ are fully faithful.
\end{lemma}

\begin{proof}
Existence of the adjoints is clear from general nonsense.
The explicit description can be seen by writing down the universal properties:
if $X$ is a cubeset and $S$ a set, then a map $S\to X(0)$ gives rise to a unique morphism
$i_!S\to X$ which factors pointwise as $i_!S(n) = S \to X(0) \to X(n)$ via the unique map $n\to 0$.
In particular, we have that evaluation at $0$ induces a natural bijection
\[
    \hom_{\cSet}(i_!S,X) \to \hom_{\set}(S,X(0)).
\]
On the other hand, a map $X(0)\to S$ induces a unique map $X(n)\to X(0)^{2^n}\to S^{2^n}$, so
\[
    \hom_{\cSet}(X,S^{2^{(-)}}) = \hom_{\set}(X(0),S).
\]
Alternatively, by the (co)limit formula for left and right Kan extension it is clear that
\begin{align*}
    i_!S(n) &= \varinjlim_{n\to 0} S = S, \\
    i_\ast S(n) &= \varprojlim_{0\to n} S = \prod_{0\to n} S = S^{2^n}.
\end{align*}
The fully faithfulness of $i_\ast$ follows directly from the fact that the counit $i^\ast i_\ast S\to S$ is an isomorphism:
$i^\ast i_\ast S$ simply is $i_\ast S(\ast) = S$.
But now the fully faithfulness of $i_*$ by general nonsense also implies fully faithfulness of $i_!$.
\end{proof}

\begin{definition}
If $S$ is a set, we call $i_! S$ the \emph{trivial cubeset} and $i_* S$ the \emph{discrete cubeset}.
\end{definition}

\begin{remark}\label{rem:left-adjoint-i_!}
The functor $i_!$ has a further left adjoint:
since $i_!$ assigns to a set $S$ the constant diagram,
its left adjoint assigns to a presheaf $X\in\cSet$ the colimit of $X$ over ${\bbox}^\op$.
In particular, $i_!$ also preserves limits.
\end{remark}

\begin{remark}
We could also consider many more triples of adjunctions between other presheaf categories, e.g. simplicial sets, by extending the diagram of index-categories
\begin{center}
\begin{tikzcd}
 \inert & {\mathrm{Fin}_{*,\le 1}^{\op}} & \bbGamma & \bbox & {\int_{\bbGamma}*/\mathbb{F}_2^{\bullet}} \\
	&& \bbDelta
	\arrow[from=1-1, to=1-2]
	\arrow[from=1-2, to=1-3]
	\arrow["j", from=1-3, to=1-4]
	\arrow["{=}"{marking, allow upside down}, draw=none, from=1-4, to=1-5]
	\arrow["{{\mathrm{Cut}}}", from=2-3, to=1-3]
\end{tikzcd}
\end{center}
where all functors are bijective on objects, the horizontal arrows are faithful and none is full,
and $\mathrm{Un}\circ j = \id$.
\end{remark}

For the rest of this subsection we will give another description of $\cSet = \PSh(\bbox)$ in terms of its unstraightening over $\bbGamma$.
Since $\bbox$ is an unstraightening of $\bbGamma$, we expect to be able to translate $\cSet$ as $\bbGamma$-sets with a certain action.
Let us explain this in more detail.

\begin{construction}
Let $\mcC$ be a small category, $G\colon\mcC\to\Grp$ a functor and let $\mathrm{Un}\colon\mcD\to\mcC$ be its unstraightening.
Given a presheaf $F\in\PSh(\mcC)$, one obtains a bifunctor $\Hom(F,F)\colon\mcC\times\mcC^\op\to\set$ 
given by precomposing the Hom-functor $\hom\colon\set^\op\times\set\to\set$ with the functor
\[
    F^\op\times F\colon\mcC\times\mcC^\op\to\set^\op\times\set.
\]
Explicitly,
\[
    \hom(F,F)(X,Y) = \hom_\set(F(X),F(Y)),\qquad X,Y\in\mcC.
\]
Moreover, we can extend $G$ to be a functor $\mcC\times\mcC^\op\to\set$ which only depends on the first
variable and uses that every group has an underlying set.

Now construct the following category $\mcE_{G,\cC}(2)$:
\begin{itemize}
    \item its objects are tuples $(F,\alpha)$ where $F$ is a presheaf on $\mcC$
        and $\alpha$ is a dinatural transformation $\alpha\colon G\to\hom(F,F)$
        whose components $\alpha_X$ are monoid homomorphisms.
	\item A morphism $\eta\colon (F,\alpha)\to(F',\alpha')$ is a natural transformation $\eta\colon F\to F'$
        such that for all $X\in\mcC$, the diagram
\begin{center}
\begin{tikzcd}
	{G(X,X)} & {\hom(F,F)(X,X)} \\
	{\hom(F',F')(X,X)} & {\hom(F(X),F'(X))}
	\arrow["{\alpha_X}", from=1-1, to=1-2]
	\arrow["{\alpha'_X}"', from=1-1, to=2-1]
	\arrow["{\eta_{X\ast}}", from=1-2, to=2-2]
	\arrow["{\eta_X^\ast}"', from=2-1, to=2-2]
\end{tikzcd}
\end{center}
is commutative.
\end{itemize}
Unwinding the definitions, the datum of the objects of $\mcE_{G,\cC}(2)$ is the same
as giving the datum of a presheaf $F\colon\mcC^\op\to\set$ together with a
$G(X)$-action $\alpha_X$ on $F(X)$ for every $X\in\mcC$ such that for each map $f\colon X\to Y$ in $\mcC$ and every $g\in G(X)$,
the diagram
\begin{center}
\begin{tikzcd}
	{F(Y)} & {F(X)} \\
	{F(Y)} & {F(X)}
	\arrow["{F(f)}", from=1-1, to=1-2]
	\arrow["{\alpha_Y(G(f)g)}"', from=1-1, to=2-1]
	\arrow["{\alpha_X(g)}", from=1-2, to=2-2]
	\arrow["{F(f)}"', from=2-1, to=2-2]
\end{tikzcd}
\end{center}
is commutative.
\end{construction}

But this information suffices to reconstruct a presheaf on $\mcD$, so we obtain the following Proposition.

\begin{proposition}\label{prop:unstrait-presehaf}
The categories $\PSh(\mcD)$ and $\mcE_{G,\cC}(2)$ are equivalent.
\end{proposition}

\begin{proof}
Restricting a presheaf $F$ on $\mcD$ to $\mcC$
yields a presheaf $F$ on $\mcC$ together with a $G(X)$-action $\alpha_X$ on $F(X)$
(which is just a monoid homomorphism $G(X)=G(X,X)\to \hom(F,F)(X,X) = \End F(X)$) by defining $\alpha_X(g)=F((\id_X, g^{-1}))$.
This action has the following property:
for every map $f\colon X\to Y$ in $\mcC$ and every $g\in G(X)$, we have that
$F(f)\circ\alpha_Y(G(f)g) = \alpha_X(g)\circ F(f)$.
This is due to the fact that in $\mcD$ we have that a morphism $(f,G(f)g^{-1})\colon X\to Y$ in $\mcD$
can be factored as $(f,e)\circ (\id_X,g^{-1})$ and as $(\id_Y,G(f)g^{-1})\circ (f,e)$,
and applying $F$ to both factorizations yields the desired equality.

In the other direction one can easily extend an object $(F,\alpha)$ in $\mcE_{G,\cC}(2)$ to a presheaf on $\mcD$
which is given on objects by $F(X)$ and a map $(f,g) = (\id_Y,g)\circ (f,e)$ in $\mcD$
is assigned to $F(f)\circ \alpha_Y(g^{-1})\colon F(Y)\to F(X)$.

Then it is routine to check that these constructions induce an equivalence of categories.
\end{proof}

This allows one to see cubesets as $\bbGamma$-sets with a compatible family of $\mathbb{F}_2^n$-actions on the set of $n$-cubes.
We can also encode the functoriality differently by not requiring a condition on the action to be compatible with some given presheaf transition maps
but rather by first giving the actions and then requiring the transition maps to be action preserving homomorphisms.

\begin{remark}
Another description of the presheaf category on $\mcD$ is the following:
define a category whose objects are tuples
$(F_X)_{X\in\mcC}$ where $F_X$ is a presheaf on $*/G(X)$ (i.e., a $G$-set/$G$-module/a set with $G$-action) together with a functorial assignment of morphisms
$\phi\colon G(f)^\ast F_Y\to F_X$ for every map $f\colon X\to Y$ in $\mcC$.
And a morphism between two such objects is a family of natural transformations,
compatible with the maps $\phi$.
Then this category also is equivalent to $\PSh(\mcD)$.
\end{remark}


\subsection{Concrete Cubesets}\label{ssec:concrete-cubesets}

The notion of concrete sheaf \cite{Dubuc1979} specializes to the topos of cubesets,
giving us the notion of concrete cubeset.
This recovers the classical notion of nilspace, since all nilspaces considered in the literature
are concrete cubesets, cf. \cite[Definition 1.2.1]{Candela2017} and \cite[Definition 3.2]{Gutman2020}.

\begin{definition}\label{def:concrete-cubeset}
A cubeset $X$ is called \emph{concrete} if the canonical map
\[
    \hom({\bbox}^n,X) = X(n) \to \prod_{0\to n} X(0) = (X(0))^{2^n}
\]
is injective.
We denote the full subcategory spanned by concrete cubesets by $\ccSet$.
\end{definition}

\begin{example}
If $S$ is a set, the trivial cubeset $i_! S$ and the discrete cubeset $i_* S$ are concrete.
Moreover, the Host-Kra cubegroups from Example~\ref{ex:host_kra_cubegroups} are concrete.
\end{example}

\begin{remark}
In view of the triple of adjoint functors from Lemma~\ref{lem:point-cset},
$X$ is concrete precisely if the unit $X\to i_\ast i^\ast X$ is a monomorphism
(recall that $i_{*}$ assigns to a set the discrete cubeset $n\mapsto X(0)^{2^n}$, and $i^*$ is simply evaluating at $0$).
Since being a monomorphism can be checked pointwise on sections, this means that the map
\[
    \hom({\bbox}^n,X)\to\hom({\bbox}^n,i_\ast i^\ast X) = \hom(i_! i^\ast {\bbox}^n,X),
\]
given by pullback along the counit $i_!i^\ast{\bbox}^n\to{\bbox}^n$ is a monomorphism.
Taking
\[
    M = \{i_!i^\ast{\bbox}^n\to{\bbox}^n : n\in\N_0\},
\]
we see that the concrete cubesets are precisely the $M$-separated presheaves on ${\bbox}^n$.
The $M$-local presheaves (=$M$-sheaves) are precisely the ones where $X \simeq i_\ast i^\ast X$,
And by fully faithfulness of $i_\ast$ this implies that the category of $M$-local presheaves is equivalent to $\Set$ under the inclusion $i_\ast\colon\set\to\cSet$.

Put differently, the geometric morphism $i\colon\set\to\cSet$ induces a Lawvere-Tierney topology on $\cSet$,
which corresponds uniquely to a Grothendieck topology on $\bbox$,
whose covers are given by the maps $i_!i^\ast{\bbox}^n\to{\bbox}^n$,
such that the sheaves for this topology are precisely $\set$, and the separated presheaves are the concrete cubesets.
See also \cite[Chapter V]{MacLane1994} and \cite[Chapter A4.3]{Johnstone2002}.
\end{remark}

\begin{lemma}\label{lem:concrete-morphism}
A map $f\colon X\to Y$ of concrete cubesets is fully determined by its value $f_0$ on points.
\end{lemma}

\begin{proof}
This follows from the commutativity of the diagram
\begin{center}
\begin{tikzcd}
	{X(n)} & {Y(n)} \\
	{\prod\limits_{0\to n}X(0)} & {\prod\limits_{0\to n}Y(0)}
	\arrow["{f_n}", from=1-1, to=1-2]
	\arrow[hook, from=1-1, to=2-1]
	\arrow[hook, from=1-2, to=2-2]
	\arrow["{\prod f_0}"', from=2-1, to=2-2]
\end{tikzcd}
\end{center}
and since the right vertical map is monic.
\end{proof}

\begin{proposition}\label{prop:concrete-left-adjoint}
The inclusion $\ccSet\to\cSet$ preserves limits and has a left adjoint $(-)^\mathrm{conc}\colon\cSet\to\ccSet$.
This left adjoint, the \emph{concretification},%
\footnote{also called concretization}
takes a cubeset $X$ and assigns to it, at stage $n$, the image of the map $X(n)\to\prod_{0\to n}X(0)$.
\end{proposition}

\begin{proof}
Being concrete is stable under limits and so $\ccSet\sub\cSet$ preserves them.
Since the image is functorial (in any elementary topos, at least),
the described construction indeed makes a concrete cubeset out of a (general) cubeset.
Given a map $f\colon X\to Y$ where $Y$ is concrete, by functoriality of the image
there exists a unique map $g_n\colon\im_n\to Y_n$ making the diagram
\begin{center}
\begin{tikzcd}
	{X(n)} & {\im_n} & {\prod\limits_{0\to n}X(0)} \\
	{Y(n)} & {Y(n)} & {\prod\limits_{0\to n}Y(0)}
	\arrow[two heads, from=1-1, to=1-2]
	\arrow["{{f_n}}", from=1-1, to=2-1]
	\arrow[hook, from=1-2, to=1-3]
	\arrow["{{g_n}}", dashed, from=1-2, to=2-2]
	\arrow["{{\prod f_0}}", from=1-3, to=2-3]
	\arrow["{=}"', from=2-1, to=2-2]
	\arrow[hook, from=2-2, to=2-3]
\end{tikzcd}
\end{center}
commutative.
This shows that taking the image is left adjoint to the inclusion $\ccSet\sub\cSet$.
\end{proof}

Differently than with limits, not every colimit of concrete cubesets is concrete.

\begin{example}\label{ex:concrete-colimit-not}
Consider the 2-corner ${\bbcor}^2\sub{\bbox}^2$ which is defined as the pushout of the inert map $\ast\to{\bbox}^1$ along itself.
Then ${\bbcor}^2$, ${\bbox}^2$ and $\ast$ are concrete cubesets, but the pushout
\begin{center}
\begin{tikzcd}
	{{\bbcor}^2} & \ast \\
	{{\bbox}^2} & {\bbox}^2/{\bbcor}^2
	\arrow[from=1-1, to=1-2]
	\arrow[from=1-1, to=2-1]
	\arrow[from=1-2, to=2-2]
	\arrow[from=2-1, to=2-2]
	\arrow["\lrcorner"{anchor=center, pos=0.125, rotate=180}, draw=none, from=2-2, to=1-1]
\end{tikzcd}
\end{center}
is not.
First of all, ${\bbox}^2/{\bbcor}^2(0) = \{0,1\}$ where $0$ corresponds to the corner ${\bbcor}^2$ and $1$ to the other point in ${\bbox}^2$.
Second, note that ${\bbox}^2/{\bbcor}^2(1)$ is the set of all maps ${\bbox}^1\to{\bbox}^2$ where two maps get identified if and only if they both factor over the corner ${\bbcor}^2\sub{\bbox}^2$.
This follows from the fact that pushouts are computed pointwise in $\cSet$ and then one checks the universal property directly.
This means that ${\bbox}^2/{\bbcor}^2(1)$ has more than four elements and so cannot be concrete since ${\bbox}^2/{\bbcor}^2(0)^2 = \{0,1\}^2$.
More explicitly, in ${\bbox}^2/{\bbcor}^2$ there exists a nontrivial 1-cube with vertices $0$ and $0$, induced by the off-diagonal.
\end{example}

However, as usual, concrete cubesets are stable under filtered colimits and coproducts.

\begin{proposition}\label{prop:concrete-stability}
Coproducts and filtered colimits of concrete cubesets are concrete.
Moreover, subobjects of concrete cubesets are concrete.
\end{proposition}

\begin{proof}
The claim about filtered colimits follows from the fact that filtered colimits commute with finite limits in $\set$.
For the claim about coproducts note that the map $\coprod_i X_i(n) \to \prod_{0\to n}\coprod_i X_i(0)$ factors as
\[
    \coprod_i X_i(n) \to \coprod_i\prod_{0\to n} X_i(0) \to \prod_{0\to n}\coprod_i X_i(0)
\]
where the second map is the natural map arising from swapping the colimit with the limit,
which in our case is injective.
The first map is just a coproduct of injections and hence also injective.

For subobjects consider an inclusion $X\to Y$ with $Y$ concrete.
Now, since monomorphisms are injections sectionwise,
by naturality of the map we obtain the commutative square
\begin{center}
\begin{tikzcd}[ampersand replacement=\&]
	{Y(n)} \& {\prod_{0\to n}Y(0)} \\
	{X(n)} \& {\prod_{0\to n}X(0)}
	\arrow[hook, from=1-1, to=1-2]
	\arrow[hook, from=2-1, to=1-1]
	\arrow[from=2-1, to=2-2]
	\arrow[from=2-2, to=1-2]
\end{tikzcd}
\end{center}
where the left vertical and the upper horizontal maps are monomorphisms.
Thus, $X(n)\to \prod_{0\to n}X(0)$ is monic as well.
\end{proof}

Moreover, concrete cubesets are stable under taking the internal Hom,
essentially due to the fact that ${\bbox}^n$ is concrete and a morphism between concrete cubesets is determined on points.

\begin{proposition}\label{prop:concrete-ihom-prod}
Let $X$ and $Y$ be cubesets, where $Y$ is concrete.
Then $\ihom(X,Y)$ is concrete.
Furthermore, concretification preserves finite products.
\end{proposition}

\begin{proof}
First we treat the case where $X$ and $Y$ are concrete, and hence so is $X\times{\bbox}^n$.
Therefore, $\hom(X\times{\bbox}^n,Y)$ consists of all maps $f_0\colon X(0)\times\{0,1\}^n\to Y(0)$ that satisfy some extra condition (see Lemma~\ref{lem:concrete-morphism}).
But $\prod_{0\to n} \hom(X,Y)$ consists of $2^n$ maps $f_0\colon X(0)\to Y(0)$ that satisfy the same kind of condition.
Therefore, the canonical map
\[
    \ihom(X,Y)(n) = \hom(X\times{\bbox}^n,Y) \to \prod_{0\to n}\hom(X,Y) = \prod_{0\to n}\ihom(X,Y)(0)
\]
is injective.
Now let $X$ be an arbitrary cubeset, and take $Z$ to be concrete.
Then we have that
\[
    \hom((X\times{\bbox}^n)^\mathrm{conc},Z) = \hom(X\times{\bbox}^n,Z) = \hom(X,\ihom({\bbox}^n,Z)).
\]
By the previous case, $\ihom({\bbox}^n,Z)$ is concrete.
Therefore, the right hand side is equal to
\[
    \hom(X,\ihom({\bbox}^n,Z)) = \hom(X^\mathrm{conc},\ihom({\bbox}^n,Z)) = \hom(X^\mathrm{conc}\times{\bbox}^n,Z) = \hom(X^\mathrm{conc}\times({\bbox}^n)^\mathrm{conc},Z)
\]
where the last equality follows from the fact that ${\bbox}^n$ is concrete.
By the Yoneda lemma we obtain that $(X\times{\bbox}^n)^\mathrm{conc} = X^\mathrm{conc}\times({\bbox}^n)^\mathrm{conc}$.
This assertion now implies that the concretification-adjunction also holds internally, i.e.,
for general $X$ and concrete $Y$ we have that
\begin{align*}
    \ihom(X,Y)(n) &= \hom(X\times{\bbox}^n,Y) \\
    &= \hom((X\times{\bbox}^n)^\mathrm{conc},Y) \\
    &= \hom(X^\mathrm{conc}\times({\bbox}^n)^\mathrm{conc},Y) \\
    &= \hom(X^\mathrm{conc}\times{\bbox}^n,Y) \\
    &= \ihom(X^\mathrm{conc},Y)(n),
\end{align*}
natural in $n$.
Thus the claim on $\ihom(X,Y)$ follows for general $X$ and concrete $Y$.
Lastly, we can deduce that $(X\times Y)^\mathrm{conc} = X^\mathrm{conc}\times Y^\mathrm{conc}$ for all cubesets $X,Y$:
If $Z$ is concrete, we have that
\begin{align*}
    \hom((X\times Y)^\mathrm{conc},Z) &= \hom(X\times Y,Z) \\
    &= \hom(X,\ihom(Y,Z)) \\
    &= \hom(X^\mathrm{conc},\ihom(Y^\mathrm{conc},Z)) \\
    &= \hom(X^\mathrm{conc}\times Y^\mathrm{conc},Z).
\end{align*}
Again, the Yoneda lemma implies that $(X\times Y)^\mathrm{conc} = X^\mathrm{conc}\times Y^\mathrm{conc}$.
Lastly, the terminal object $\ast$ is concrete, so $\ast^\mathrm{conc}=\ast$.
\end{proof}

\begin{remark}
Note that the inclusion $\ccSet\sub\cSet$ is not a geometric morphism
since concretification doesn't preserve all finite limits.

For example, consider a non-surjective map $X\to Y$  between cubesets $X$ and $Y$ inducing a bijection on underlying sets $X(0)\simeq Y(0)$, where $Y$ is concrete
(e.g., consider the counit $i_! i^* X\to X$ for some non-constant concrete $X$).
Now build the quotient $Y/X$ (the cofiber) so that the diagram
\begin{center}
\begin{tikzcd}[ampersand replacement=\&]
	X \& Y \\
	{*} \& {Y/X}
	\arrow[from=1-1, to=1-2]
	\arrow[from=1-1, to=2-1]
	\arrow[from=1-2, to=2-2]
	\arrow[from=2-1, to=2-2]
\end{tikzcd}
\end{center}
is a pushout square.
Consider the kernel pair $K$ of the map $Y\to Y/X$.
Then $K$ is a subobject of the concrete cubeset $Y\times Y$ and hence concrete.
Since $X\ne Y$, we have $Y/X\ne \ast$ and $Y\to Y/X$ is a surjection.
Thus, $K$ is a proper subobject of $Y\times Y$.
On the other hand, the concretification of $Y/X$ is the point $\ast$ since $Y/X(\ast)=\ast$.
So we conclude
\[
    (Y\times_{Y/X}Y)^\mathrm{conc} = K^\mathrm{conc} = K \neq Y\times Y = Y^\mathrm{conc}\times_{(Y/X)^\mathrm{conc}}Y^\mathrm{conc}.
\]
\end{remark}
